% Vietnamese translation for OI Public Library
% Source-PDF: IOI中国国家候选队论文/国家集训队2023论文集.pdf
\documentclass[11pt]{article}
\usepackage{oipl-vn}

\title{Tập luận văn đội tuyển quốc gia Trung Quốc năm 2023}
\author{Huấn luyện viên: Yang Yaoliang, Peng Sijin}
\date{China Computer Federation, 2023}

\begin{document}
\maketitle

\origtitle{Tập luận văn đội tuyển quốc gia Trung Quốc Olympic Tin học năm 2023}
\sourcepath{IOI China candidate team papers / National training team 2023 collection.pdf}

\renewcommand{\abstractname}{Tóm tắt}
\begin{abstract}
PDF nguồn là tập hợp luận văn đội tuyển quốc gia Trung Quốc năm 2023. Tệp được
tạo bằng LaTeX với lớp văn bản Unicode đọc được; mục lục và các trang thân bài
đã kiểm tra đều trích xuất tốt bằng \texttt{pdftotext}. Bản LaTeX này chuyển
ngữ phần nhận diện tập sách và toàn bộ mục lục, đồng thời ghi lại các chủ đề
chính để phục vụ bản dịch chi tiết hơn.
\end{abstract}

\section{Thông tin đầu sách}

Tập sách có tiêu đề \emph{Tập luận văn đội tuyển quốc gia Trung Quốc Olympic
Tin học năm 2023}. Trang bìa ghi huấn luyện viên là Yang Yaoliang và Peng
Sijin. Tệp PDF gồm 249 trang và được tạo bằng LaTeX với gói \texttt{hyperref}.

\section{Mục lục đã dịch}

\begin{center}
\small
\begin{tabular}{r p{0.52\linewidth} p{0.24\linewidth}}
\textbf{Trang} & \textbf{Bài viết} & \textbf{Tác giả} \\
\hline
1 & Tổng quan một số phương pháp xử lý tính liên thông và các bài toán liên quan trong lý thuyết đồ thị & Wan Chengzhang \\
31 & Bàn về xây dựng và sinh dữ liệu trong thi tin học & Liu Yiping \\
50 & Bàn về một lớp bài toán có treewidth bị chặn & Wu Chang \\
63 & Bao hàm--loại trừ trên quan hệ đôi một khác nhau & Zhou Ziheng \\
70 & Báo cáo ra đề \emph{Nhà ảo thuật} & Meng Yuhao \\
85 & Hình học trong OI & Chang Ruinian \\
104 & Bàn về các thuật toán liên quan tới xâu palindrome và ứng dụng & Xu Anyi \\
113 & Bàn về một số ứng dụng của trường hữu hạn trong OI & Qi Langrui \\
123 & Bàn về một lớp bài toán thống kê trên cây & Zhu Yikai \\
131 & Bàn về ứng dụng của ma trận trong thi tin học & Yang Ningyuan \\
147 & Bàn về một số bài toán liên quan tới matching đồ thị hai phía & Ke Yisi \\
156 & Bàn về ứng dụng đặc biệt của DSU trong OI & Wang Xiangwen \\
167 & Bàn về ứng dụng static top tree trên cây và đồ thị nối tiếp-song song tổng quát & Cheng Siyuan \\
181 & Bước đầu về hypercube và các thuật toán liên quan & Guan Yanru \\
198 & Bàn về một số phương pháp phân tích thừa số nguyên tố & Luo Siyuan \\
209 & Một lớp cấu trúc dữ liệu xâu con cơ bản & Xu Tingqiang \\
220 & Xem xét lại một vài bài toán số học cổ điển & Zheng Xuanye \\
232 & Bàn về ứng dụng tính lồi của hàm số trong OI & Guo Yuchong
\end{tabular}
\end{center}

\section{Chủ đề chính}

Từ mục lục và phần đầu của bài thứ nhất, tập năm 2023 bao phủ các nhóm chủ đề
sau:
\begin{itemize}
  \item tính liên thông trong đồ thị, treewidth, đồ thị nối tiếp-song song và static top tree;
  \item xây dựng dữ liệu, bao hàm--loại trừ và các báo cáo ra đề;
  \item hình học, trường hữu hạn, ma trận, tính lồi và số học;
  \item palindrome, cấu trúc dữ liệu xâu con và DSU;
  \item matching đồ thị hai phía, hypercube và phân tích thừa số nguyên tố;
  \item thống kê trên cây và các bài toán cổ điển được xem xét lại.
\end{itemize}

\section{Ghi chú về trích xuất}

Phần thân của tập 2023 có thể trích xuất được mà không cần OCR. Bài đầu tiên
chia nội dung thành đồ thị vô hướng và đồ thị có hướng, trình bày các ý tưởng
thuật toán về tính liên thông trong thi tin học, đồng thời bàn thêm các bài
toán đồ thị có liên hệ với tính liên thông. Tập này có chất lượng văn bản phù
hợp để mở rộng thành bản dịch từng bài với công thức, hình và bảng.

\end{document}
